%! Composition of Functions — Unit 1, Lesson 1.5 — Slides (Beamer)
\documentclass[aspectratio=169, 11pt]{beamer}
\usepackage{precalculus-beamer}   % provides \plumheader{} and \sectionlabel[color]{}

\begin{document}

% TITLE SLIDE
{
\setbeamercolor{background canvas}{bg=plum}
\begin{frame}[plain]
  \vspace{2.8cm}\hspace{0.55cm}%
  \begin{minipage}{0.9\textwidth}
    {\color{goldacc}\bfseries\small LESSON 1.5}\\[0.5cm]
    {\color{white}\bfseries\LARGE Composition of Functions}\\[1.2cm]
    {\color{lilacmid}\small Precalculus~$\cdot$~Unit 1: Functions and Change}
  \end{minipage}
\end{frame}
}

% HOOK SLIDE
\begin{frame}[t, plain]
  \plumheader{Hook: Two Registers, One Jacket}
  \sectionlabel[goldacc]{Think About It}

  A \$60 jacket. A \$10 coupon. Six percent sales tax.

  \bigskip
  \begin{columns}[T]
    \column{0.46\textwidth}
      \begin{block}{Register A --- coupon first}
        Take \$10 off, \textit{then} add the tax.\\[4pt]
        \(1.06(60 - 10) = \$53.00\)
      \end{block}
    \column{0.46\textwidth}
      \begin{block}{Register B --- tax first}
        Add the tax, \textit{then} take \$10 off.\\[4pt]
        \(1.06(60) - 10 = \$53.60\)
      \end{block}
  \end{columns}

  \bigskip
  \textit{Same jacket, same coupon, same tax rate --- sixty cents apart. Today
  we chain functions together, and the order is part of the answer.}
\end{frame}

% SECTION 1 — TWO MACHINES IN A ROW
\begin{frame}[t, plain]
  \plumheader{1. Two Machines in a Row}

  \begin{center}
    {\Large \((f \circ g)(x) = f\big(g(x)\big)\)} \quad --- read ``\(f\) of
    \(g\) of \(x\).''
  \end{center}

  \bigskip
  \begin{center}
  \begin{tikzpicture}[x=1cm, y=1cm]
    \draw[thick, plum, fill=lilac, rounded corners=2pt]
      (-0.85,-0.55) rectangle (0.85,0.55);
    \node at (0,0) {\large \(g\)};
    \draw[thick, plum, fill=lilac, rounded corners=2pt]
      (4.15,-0.55) rectangle (5.85,0.55);
    \node at (5,0) {\large \(f\)};
    \draw[->, very thick, plum] (-2.6,0) -- (-0.95,0);
    \node[above, font=\small] at (-1.75,0.06) {\(x\)};
    \draw[->, very thick, plum] (0.95,0) -- (4.05,0);
    \node[above, font=\small] at (2.5,0.06) {\(g(x)\)};
    \draw[->, very thick, plum] (5.95,0) -- (8.2,0);
    \node[above, font=\small] at (7.05,0.06) {\(f(g(x))\)};
    \node[below, font=\small\itshape, text=slate] at (0,-0.62) {runs first};
    \node[below, font=\small\itshape, text=slate] at (5,-0.62) {runs second};
  \end{tikzpicture}
  \end{center}

  \bigskip
  \begin{block}{The one sentence for today}
    Work from the \textbf{inside out}. The function written \textbf{last} is
    the one that runs \textbf{first} --- and the circle is \textit{not}
    multiplication.
  \end{block}
\end{frame}

% SECTION 2 — VALUES FROM A TABLE
\begin{frame}[t, plain]
  \plumheader{2. Composite Values from a Table}

  \begin{columns}[T]
    \column{0.44\textwidth}
      \begin{tabular}{|c|c|c|c|c|c|}
      \hline
      \(x\) & 0 & 1 & 2 & 3 & 4 \\
      \hline
      \(f(x)\) & 3 & 4 & 0 & 2 & 1 \\
      \hline
      \(g(x)\) & 2 & 3 & 7 & 0 & 4 \\
      \hline
      \end{tabular}

      \bigskip
      \textit{Write the middle number down. Every time.}

    \column{0.52\textwidth}
      \begin{block}{\(f(g(1))\)}
        \(g(1) = 3\), \quad then \(f(3) = \mathbf{2}\)
      \end{block}

      \medskip
      \begin{block}{\(g(f(1))\)}
        \(f(1) = 4\), \quad then \(g(4) = \mathbf{4}\)
      \end{block}

      \medskip
      \begin{block}{\(f(g(2))\) --- careful}
        \(g(2) = 7\), and \(f(7)\) is not in the table.\\[3pt]
        The inner output has to be something \(f\) accepts.
      \end{block}
  \end{columns}
\end{frame}

% SECTION 3 — BUILDING ANALYTICALLY
\begin{frame}[t, plain]
  \plumheader{3. Building a Composite Analytically}

  \begin{center}
    To build \(f(g(x))\), substitute \(g(x)\) for \textbf{every} \(x\) in
    \(f(x)\).
  \end{center}

  \bigskip
  Let \(f(x) = 2x + 1\) and \(g(x) = x^2 - 3\).

  \bigskip
  \begin{columns}[T]
    \column{0.46\textwidth}
      \begin{block}{\(f(g(x))\)}
        \(f(x^2 - 3) = 2(x^2 - 3) + 1\)\\[3pt]
        \(= \mathbf{2x^2 - 5}\)
      \end{block}
    \column{0.46\textwidth}
      \begin{block}{\(g(f(x))\)}
        \(g(2x + 1) = (2x + 1)^2 - 3\)\\[3pt]
        \(= \mathbf{4x^2 + 4x - 2}\)
      \end{block}
  \end{columns}

  \bigskip
  \textit{Not even the same shape of expression. You already did this move in
  the warm-up: \(f(x^2) = 2x^2 + 1\).}
\end{frame}

% SECTION 4 — ORDER MATTERS
\begin{frame}[t, plain]
  \plumheader{4. Order Matters --- the Checkout Counter}

  \begin{center}
    \textbf{Coupon:} \(c(p) = p - 10\) \qquad
    \textbf{Tax:} \(t(x) = 1.06x\)
  \end{center}

  \bigskip
  \begin{columns}[T]
    \column{0.46\textwidth}
      \begin{block}{Coupon first: \(t \circ c\)}
        \(t(c(p)) = 1.06(p - 10)\)\\[3pt]
        \(= 1.06p - 10.60\)\\[3pt]
        At \(p = 60\): \(\$53.00\)
      \end{block}
    \column{0.46\textwidth}
      \begin{block}{Tax first: \(c \circ t\)}
        \(c(t(p)) = 1.06p - 10\)\\[3pt]
        \phantom{\(= 1.06p - 10.60\)}\\[3pt]
        At \(p = 60\): \(\$53.60\)
      \end{block}
  \end{columns}

  \bigskip
  \begin{block}{Why exactly sixty cents, at every price?}
    Taxing first charges 6\% on the \$10 you never paid: \(0.06(10) = \$0.60\).
    Neither function alone answers ``what do I pay?'' --- the \textbf{chain}
    does.
  \end{block}
\end{frame}

% SECTION 5 — DECOMPOSITION, IDENTITY, LESSON 1.4
\begin{frame}[t, plain]
  \plumheader{5. Decomposition, Identity, and Lesson 1.4}

  \begin{columns}[T]
    \column{0.46\textwidth}
      \begin{block}{Decompose \(h(x) = (3x+2)^4\)}
        With a calculator, what would you do \textbf{first}? That is the inner
        function.\\[4pt]
        Inner \(g(x) = 3x + 2\)\\[3pt]
        Outer \(f(u) = u^4\)
      \end{block}

      \medskip
      \begin{block}{Identity: \(I(x) = x\)}
        \(I(f(x)) = f(x)\) and \(f(I(x)) = f(x)\).\\[3pt]
        Remember this one --- Lesson 1.6 hunts for it.
      \end{block}

    \column{0.5\textwidth}
      \begin{block}{Lesson 1.4 was composition all along}
        \(f(x + 3) = f(h(x))\) with \(h(x) = x + 3\)\\[4pt]
        \(f(2x) = f(h(x))\) with \(h(x) = 2x\)\\[6pt]
        The inner function reaches the \(x\) \textbf{before} \(f\) does ---
        which is exactly why an ``inside'' change rewrites the
        \textbf{inputs}.
      \end{block}
  \end{columns}
\end{frame}

% CLASS PRACTICE
\begin{frame}[t, plain]
  \plumheader{Class Practice}
  \sectionlabel[redacc]{Try It}

  \begin{enumerate}\setlength{\itemsep}{10pt}
    \item From the table above: find \(f(g(3))\) and \(g(f(3))\). Write the
          middle number down each time.
    \item Let \(f(x) = x^2\) and \(g(x) = x - 5\). Build \((f \circ g)(x)\)
          and \((g \circ f)(x)\), simplified. Are they the same function?
    \item Decompose \(h(x) = \sqrt{4x - 7}\) into an inner and an outer
          function --- then check by substituting back.
    \item A \$45 shirt, the \$10 coupon, and 6\% tax. Which order is cheaper
          for the shopper, and by how much?
  \end{enumerate}
\end{frame}

\end{document}
